\chapter{Simulation Architecture: S11 Protein Folding Engine}
\label{ch:protein_sim_architecture}

This chapter documents the complete computational architecture of the axiom-derived protein folding engine.  Every equation is mapped to the corresponding implementation in the physics engine, and the full reproduction procedure is provided.

The engine has a single objective function:
\begin{align}
    \mathcal{L}(\mathbf{r}) &= \underbrace{\overline{|S_{11}|^2}}_{\text{impedance mismatch}} + \underbrace{\frac{\lambda_b}{N} \sum_i (d_i - d_0)^2}_{\text{bond integrity}} \notag \\
    &\quad + \underbrace{\frac{\lambda_s}{N} \sum_{|i-j|\geq 3} \max(0, r_{\text{steric}} - d_{ij})^2}_{\text{Pauli exclusion}} + \underbrace{\mathcal{L}_{\text{port}}}_{\text{cavity coupling}}
    \label{eq:total_loss}
\end{align}
where $\mathbf{r} \in \mathbb{R}^{N \times 3}$ are the $C_\alpha$ coordinates, $N$ is the chain length, and all gradients $\partial \mathcal{L}/\partial \mathbf{r}_i$ are computed by JAX automatic differentiation.

\section{Tier~1: Axiom-Derived Constants}

Every physical constant in the engine traces to the four AVE axioms through the physics engine module \texttt{ave.solvers.protein\_bond\_constants}.

\begin{table}[H]
\centering
\caption{Engine constants and their derivation chain.  No empirical structural data enters.}
\label{tab:engine_constants_protein}
\small
\begin{tabular}{@{}l c l l@{}}
\toprule
\textbf{Symbol} & \textbf{Value} & \textbf{Derivation} & \textbf{Source} \\
\midrule
$d_0$       & 3.8~\AA       & $C_\alpha$--$C_\alpha$ peptide geometry  & Axioms 1--2 $\to$ covalent radii \\
$r_{C_\alpha}$ & 1.7~\AA    & Slater orbital radius (Carbon, $n^*=2$) & Axioms 1--2 $\to$ periodic table \\
$Q$         & 7.0           & Amide-V: $f_0/\Delta f = 23/3.3$~THz    & Axiom 1 $\to$ backbone $f_\text{res}$ \\
$Z_0$       & 1.0           & Normalised backbone impedance            & Definition \\
$\kappa$    & 0.5           & Critical coupling ($Q_\text{ext} = Q_\text{int}$) & Unique resonator maximum \\
$R_\text{burial}$ & 7.6~\AA & $2 \times d_0$ (helix contact diameter) & Axiom-derived geometry \\
$D_{\mathrm{H_2O}}$ & 2.75~\AA  & Water molecular diameter                  & Axiom 2 $\to$ O--H bond length \\
$\beta_\text{burial}$ & 1.6~\AA$^{-1}$ & $4.4 / D_{\mathrm{H_2O}}$ (sigmoid 10--90\% = $D_{\mathrm{H_2O}}$) & Derived \\
$r_\text{steric}$ & 3.4~\AA & $2 \times r_{C_\alpha}$ (Pauli exclusion) & Axiom-derived Slater radii \\
$Z_{\mathrm{H_2O}}$ & $\sqrt{80} \approx 8.9$ & $\sqrt{\varepsilon_r}$ for liquid water & Axiom 1 $\to$ dielectric theory \\
$\delta_\chi$ & 0.05~rad    & Ramachandran asymmetry / $Q$             & $\frac{1}{Q} \times 0.35$ \\
$\chi_0$    & 5.0~\AA$^3$   & $d_0^3 / 11$ (helix unit cell / geometry) & Derived from $d_0$ \\
\bottomrule
\end{tabular}
\end{table}

\section{Tier~2: Complex Topological Impedance $Z_{topo}$}
\label{sec:z_topo_table}

Each amino acid maps to a complex impedance $Z_i = R_i + jX_i$ where:

\begin{itemize}
    \item $R_i$ encodes the sidechain's hydrophobic volume (resistive coupling to vacuum)
    \item $X_i$ encodes the sidechain's charge reactance, scaled by the backbone $Q$-factor:
\end{itemize}

\begin{equation}
    X_i = \begin{cases}
        -R_i / Q     & \text{negative charge (capacitive): D, E} \\
        +R_i / Q     & \text{positive charge (inductive): K, R} \\
        \pm R_i / 2Q & \text{polar uncharged: S, T, C, Y, N, Q, H} \\
        0            & \text{hydrophobic: A, V, I, L, M, F, W, P, G}
    \end{cases}
    \label{eq:x_charge}
\end{equation}

The $Q$-factor scaling ensures that resistive (hydrophobic) coupling dominates at $\sim 85\%$, with reactive (electrostatic) coupling as a $\sim 15\%$ perturbation.  This ratio is \emph{derived} from the backbone's amide-V resonance width.

\begin{table}[H]
\centering
\caption{Canonical complex $Z_{topo}$ table from \texttt{protein\_bond\_constants.py}.  $|Z|$ is the magnitude used in the ABCD cascade.  Values derive from sidechain topology via the R-group shunt impedance model (Chapter~\ref{ch:biological_circuitry}).}
\label{tab:z_topo_complex}
\small
\begin{tabular}{@{}l c c r l@{}}
\toprule
\textbf{AA} & $R$ & $X$ & $|Z|$ & \textbf{Type} \\
\midrule
G & 0.50 & 0.000 & 0.50 & Hydrophobic (minimal stub) \\
A & 0.53 & 0.000 & 0.53 & Hydrophobic \\
E & 0.52 & $-0.074$ & 0.53 & Negative charge \\
R & 0.55 & $+0.079$ & 0.56 & Positive charge \\
K & 0.60 & $+0.086$ & 0.61 & Positive charge \\
Q & 0.63 & $+0.045$ & 0.63 & Polar uncharged \\
D & 0.66 & $-0.094$ & 0.67 & Negative charge \\
I & 0.73 & 0.000 & 0.73 & Hydrophobic \\
M & 0.87 & 0.000 & 0.87 & Hydrophobic \\
V & 0.93 & 0.000 & 0.93 & Hydrophobic \\
L & 1.00 & 0.000 & 1.00 & Hydrophobic \\
N & 1.10 & $+0.079$ & 1.10 & Polar uncharged \\
Y & 1.31 & $-0.094$ & 1.31 & Polar uncharged \\
F & 1.57 & 0.000 & 1.57 & Hydrophobic \\
S & 1.64 & $+0.117$ & 1.64 & Polar uncharged \\
T & 1.73 & $+0.124$ & 1.74 & Polar uncharged \\
C & 1.74 & $-0.124$ & 1.74 & Polar uncharged \\
H & 2.50 & $+0.179$ & 2.51 & Polar uncharged (half-protonated) \\
W & 3.40 & 0.000 & 3.40 & Hydrophobic \\
P & 5.02 & 0.000 & 5.02 & Hydrophobic (cyclic lock) \\
\bottomrule
\end{tabular}
\end{table}


\section{Tier~3: Physics Layers}

The loss function (Eq.~\ref{eq:total_loss}) consists of eight axiom-derived physics layers evaluated in sequence.  Each layer contributes either to the S$_{11}$ cascade loss, the bond/steric penalty, or the port coupling loss.

\subsection{Layer 1: Conjugate Impedance Matching}

Non-local residue pairs ($|i-j| > 2$, $d_{ij} < 15$~\AA) couple through a complex admittance based on their conjugate impedance product:

\begin{equation}
    Y_{ij}^{\text{shunt}} = \frac{\kappa \cdot m_{ij} \cdot C_{ij}^{\text{sat}}}{d_{ij}^2}
    \label{eq:shunt_admittance}
\end{equation}
where the normalised conjugate match quality is:
\begin{equation}
    m_{ij} = \max\!\left(0,\; \frac{\operatorname{Re}(Z_i \, Z_j^*)}{|Z_i|\,|Z_j|}\right)
    \label{eq:match_quality}
\end{equation}

The conjugate product $\operatorname{Re}(Z_i Z_j^*) = R_i R_j + X_i X_j$ is:
\begin{itemize}
    \item \textbf{Positive} for hydrophobic pairing ($R \times R > 0$) and salt bridges ($+jX \times -jX > 0$)
    \item \textbf{Negative} for like-charge repulsion ($+jX \times +jX < 0$), clamped to zero
\end{itemize}

Thus Coulomb-like repulsion \emph{emerges} from the gradient: bringing like-charged residues close does not lower $m_{ij}$ (it stays at zero), so the gradient provides no coupling force---and the steric penalty pushes them apart.

\subsection{Layer 2: Axiom~4 Dielectric Saturation}

The capacitive saturation factor from Axiom~4 amplifies coupling between well-matched close pairs:

\begin{equation}
    C_{ij}^{\text{sat}} = 1 + \left(\frac{1}{\sqrt{1 - (d_0/d_{ij})^2}} - 1\right) \cdot m_{ij}
    \label{eq:c_sat}
\end{equation}

The ratio $d_0/d_{ij}$ is clipped to $[0, 0.95]$ for numerical stability.  When a pair is well-matched \emph{and} close ($d_{ij} \lesssim d_0$), the saturation factor can reach $C^{\text{sat}} \approx 3\times$, providing strong gradient signal for helix packing.  When mismatched, $m_{ij} = 0$ suppresses the amplification entirely.

\subsection{Layer 3: Solvent Impedance Boundary}

Every residue couples to the aqueous environment through a parasitic shunt to ground:
\begin{equation}
    Y_i^{\text{solvent}} = \frac{e_i}{Z_{\mathrm{H_2O}}}, \qquad e_i = \operatorname{clip}\!\left(1 - \frac{n_i}{n_\text{max}},\; 0,\; 1\right)
    \label{eq:solvent_shunt}
\end{equation}
where $e_i$ is the fractional exposure (1 = fully solvent-exposed, 0 = fully buried), $n_i$ is the smooth neighbour count:
\begin{equation}
    n_i = \sum_{|j-i|>2} \sigma\!\left(\beta_\text{burial}(R_\text{burial} - d_{ij})\right)
    \label{eq:smooth_neighbors}
\end{equation}
and $n_\text{max} = \max(N/3, 4)$.  This creates a penalty for exposing hydrophobic residues: their high $R_i$ stubs reflect energy back into the backbone, but when exposed, the solvent shunt drains that energy to ground.  Burial eliminates the shunt, reducing $S_{11}$.

\subsection{Layer 4: ABCD Transmission Line Cascade}

The backbone is modelled as $N-1$ cascaded transmission line sections.  Each section $i$ has:
\begin{itemize}
    \item Characteristic impedance: $Z_c^{(i)} = \frac{1}{2}(|Z_i| + |Z_{i+1}|)$
    \item Electrical length: $\beta \ell^{(i)} = \omega \cdot d_i / d_0 - \delta_\chi^{(i)}$
    \item Shunt admittance at junction: $Y_i = Y_i^{\text{shunt}} + Y_i^{\text{solvent}}$
\end{itemize}

The ABCD matrix for section $i$ is:
\begin{equation}
    \mathbf{T}_i = \underbrace{\begin{pmatrix} \cos\beta\ell & jZ_c\sin\beta\ell \\ j\sin\beta\ell/Z_c & \cos\beta\ell \end{pmatrix}}_{\text{transmission line}} \cdot \underbrace{\begin{pmatrix} 1 & 0 \\ Y_i & 1 \end{pmatrix}}_{\text{shunt element}}
    \label{eq:abcd_section}
\end{equation}

The total cascade is $\mathbf{T} = \prod_{i=0}^{N-2} \mathbf{T}_i$, computed via a JIT-compiled \texttt{jax.lax.fori\_loop} for $O(N)$ scaling with zero Python overhead.

The port reflection coefficient is:
\begin{equation}
    S_{11} = \frac{A + B/Z_0 - CZ_0 - D}{A + B/Z_0 + CZ_0 + D}
    \label{eq:s11_from_abcd}
\end{equation}

\subsection{Layer 5: Multi-Frequency Integration}

The S$_{11}$ is averaged over five frequencies spanning the backbone resonance $\pm$ harmonics:
\begin{equation}
    \overline{|S_{11}|^2} = \frac{1}{5} \sum_{k=1}^{5} |S_{11}(\omega_k)|^2, \qquad \omega_k / \omega_0 \in \{0.5,\; 0.8,\; 1.0,\; 1.3,\; 2.0\}
    \label{eq:multi_freq}
\end{equation}
This prevents the optimizer from ``gaming'' a single frequency and ensures broadband impedance matching---the physical requirement for a structure bathed in broadband thermal noise.

\subsection{Layer 6: Non-Reciprocal Chirality Phase}

The AVE vacuum lattice (SRS/K4 net) is intrinsically chiral.  This enters the ABCD cascade as a non-reciprocal phase correction:
\begin{equation}
    \delta_\chi^{(i)} = \frac{0.35}{Q} \cdot w_\text{helix}^{(i)} \cdot \tanh\!\left(\frac{\chi_i}{\chi_0}\right)
    \label{eq:chiral_phase}
\end{equation}
where:
\begin{itemize}
    \item $\chi_i = (\mathbf{b}_{i} \times \mathbf{b}_{i+1}) \cdot \mathbf{b}_{i+2}$ is the scalar triple product of consecutive bond vectors (positive for right-handed twist)
    \item $w_\text{helix}^{(i)} = \max(0, 1 - \bar{Z}^{(i)}/2)$ suppresses chirality for sheet-forming segments
    \item $\chi_0 = d_0^3 / 11 \approx 5.0$~\AA$^3$ normalises the triple product
    \item $\tanh$ saturates the signal to $[-1, 1]$
\end{itemize}

This is the sole source of L/D selectivity in the engine: right-handed helices receive a phase bonus ($\delta_\chi > 0$) that lowers their electrical length and reduces $S_{11}$.

\subsection{Layer 7: Cross-Coupled Cavity Filter}

A folded protein is a set of coupled resonant cavities, not a single cascade.  The engine identifies segment boundaries via the local reflection coefficient:
\begin{equation}
    \Gamma_j = \frac{|Z_{j+1}| - |Z_j|}{|Z_{j+1}| + |Z_j|}, \qquad \text{turn indicator: } t_j = \sigma(20(\Gamma_j - 0.3))
\end{equation}

\paragraph{Layer 7a: Adjacent Junction Coupling.}
At each detected turn, the transmission between adjacent segments is:
\begin{equation}
    S_{21}^{(j)} = (1 - \Gamma_j^2) \cdot \frac{2\bar{Z}_L \bar{Z}_R}{\bar{Z}_L^2 + \bar{Z}_R^2} \cdot \exp\!\left(-\frac{|\mathbf{c}_L - \mathbf{c}_R|}{R_\text{burial}}\right)
    \label{eq:junction_s21}
\end{equation}
where $\bar{Z}_{L,R}$ and $\mathbf{c}_{L,R}$ are the mean impedance and centroid of the left/right segments.

\paragraph{Layer 7b: Non-Adjacent Cross-Coupling.}
For all segment pairs $(p, q)$ with $|p - q| \geq 2$ (skipping the nearest neighbour):
\begin{equation}
    S_{21}^{(p,q)} = \frac{2\bar{Z}_p \bar{Z}_q}{\bar{Z}_p^2 + \bar{Z}_q^2} \cdot \exp\!\left(-\frac{|\mathbf{c}_p - \mathbf{c}_q|}{R_\text{burial}}\right)
    \label{eq:cross_s21}
\end{equation}

The port coupling loss rewards high $|S_{21}|^2$:
\begin{equation}
    \mathcal{L}_\text{port} = \underbrace{\frac{1}{N}\sum_j t_j(|S_\text{self}|^2 - |S_{21}|^2)}_{\text{junction loss}} - \underbrace{\frac{1}{N}\sum_{p<q, |p-q|\geq 2} |S_{21}^{(p,q)}|^2}_{\text{cross-coupling bonus}}
    \label{eq:port_loss}
\end{equation}

\subsection{Layer 8: Bond Integrity and Steric Exclusion}

Two hard physical constraints are added as regularisation penalties:
\begin{align}
    \mathcal{L}_\text{bond}    &= \frac{2}{N} \sum_{i=0}^{N-2} (d_{i,i+1} - d_0)^2 \label{eq:bond_penalty} \\
    \mathcal{L}_\text{steric}  &= \frac{1}{N} \sum_{\substack{i < j \\ |i-j| \geq 3}} \max(0,\; r_\text{steric} - d_{ij})^2 \label{eq:steric_penalty}
\end{align}

These are not ``forces''---they are smooth, differentiable loss terms whose gradients enforce backbone connectivity ($d_i \approx 3.8$~\AA) and prevent chain self-intersection ($d_{ij} > 3.4$~\AA).

\section{Optimizer Architecture}

\subsection{Adam with Simulated Annealing}

The coordinates $\mathbf{r}$ are optimised using the Adam optimiser (Kingma \& Ba, 2015) implemented via \texttt{optax.adam(lr=10^{-3})}:

\begin{equation}
    \mathbf{r}^{(t+1)} = \mathbf{r}^{(t)} - \eta \cdot \frac{\hat{m}_t}{\sqrt{\hat{v}_t} + \epsilon}
    \label{eq:adam}
\end{equation}

Gradients are clipped to $|\mathbf{g}| \leq 10$ to prevent divergence.  For the first 50\% of training steps, Gaussian noise with temperature schedule $T(t) = 0.02(1 - 2t/n_\text{steps})^2$ is added to explore the energy landscape:
\begin{equation}
    \mathbf{r}^{(t+1)} \mathrel{+}= \mathcal{N}(0, T(t)) \qquad \text{for } t < n_\text{steps}/2
    \label{eq:annealing}
\end{equation}

The chain is re-centred at its centre of mass after every step.

\subsection{Z-Dependent Initialisation}

The initial geometry is seeded using the $Z_{topo}$ of each residue:
\begin{itemize}
    \item $|Z_i| \leq 1.0$ (helix-forming): helical seed with $100^\circ$/residue rotation and $1.5$~\AA\ rise per turn
    \item $|Z_i| > 1.0$ (sheet-forming): extended zigzag conformation
\end{itemize}
A small Gaussian perturbation ($\sigma = 0.15$~\AA) breaks initial symmetry.

\section{Script Reproduction Procedure}

The entire computation is reproduced in two commands from the repository root:

\begin{verbatim}
# Fold three test sequences (Polyalanine, Chignolin, Trpzip2)
PYTHONPATH=src python scripts/book_5_topological_biology/s11_fold_engine_v3_jax.py

# Generate SPICE transmission line strain plot
PYTHONPATH=src python scripts/book_5_topological_biology/ \
    simulate_protein_spice_transmission_line.py
\end{verbatim}

\subsection{Code--Equation Map}

Table~\ref{tab:code_equation_map} maps every equation in this chapter to its implementation in the engine source.

\begin{table}[H]
\centering
\caption{Equation-to-code traceability.  All implementations are in \texttt{s11\_fold\_engine\_v3\_jax.py} unless noted.}
\label{tab:code_equation_map}
\small
\begin{tabular}{@{}l l l@{}}
\toprule
\textbf{Equation} & \textbf{Function / Line} & \textbf{Physics Layer} \\
\midrule
Eq.~\ref{eq:total_loss}           & \texttt{\_s11\_loss()}, L299            & Total loss function \\
Eq.~\ref{eq:match_quality}        & \texttt{conjugate\_match}, L118--120    & Conjugate impedance \\
Eq.~\ref{eq:c_sat}                & \texttt{C\_sat}, L133--135              & Axiom 4 saturation \\
Eq.~\ref{eq:shunt_admittance}     & \texttt{coupling}, L137                 & Shunt admittance \\
Eq.~\ref{eq:solvent_shunt}        & \texttt{Y\_solvent}, L147--154          & Solvent boundary \\
Eq.~\ref{eq:abcd_section}         & \texttt{cascade\_step()}, L196--213     & ABCD cascade \\
Eq.~\ref{eq:s11_from_abcd}        & \texttt{s11\_at\_freq()}, L218--221     & $S_{11}$ extraction \\
Eq.~\ref{eq:multi_freq}           & \texttt{s11\_total} loop, L224--227     & Multi-frequency \\
Eq.~\ref{eq:chiral_phase}         & \texttt{chiral\_correction}, L161--183  & Chirality phase \\
Eq.~\ref{eq:junction_s21}         & \texttt{junction\_s21()}, L244--259     & Adjacent coupling \\
Eq.~\ref{eq:cross_s21}            & \texttt{cross\_loss} loop, L266--284    & Cross-coupling \\
Eq.~\ref{eq:bond_penalty}         & \texttt{bond\_penalty}, L288--290       & Bond integrity \\
Eq.~\ref{eq:steric_penalty}       & \texttt{steric\_penalty}, L292--297     & Pauli exclusion \\
Eq.~\ref{eq:adam}                  & \texttt{fold\_s11\_jax()}, L358--367    & Adam optimiser \\
Eq.~\ref{eq:annealing}            & \texttt{anneal} block, L369--374        & Simulated annealing \\
\bottomrule
\end{tabular}
\end{table}

\subsection{Dependencies}

\begin{table}[H]
\centering
\caption{External dependencies for the protein folding engine.}
\label{tab:dependencies}
\small
\begin{tabular}{@{}l l l@{}}
\toprule
\textbf{Package} & \textbf{Version} & \textbf{Purpose} \\
\midrule
\texttt{jax}   & $\geq 0.4$ & Automatic differentiation, JIT compilation \\
\texttt{optax} & $\geq 0.2$ & Adam optimiser implementation \\
\texttt{numpy} & $\geq 1.24$ & Array manipulation, output analysis \\
\texttt{control} & $\geq 0.9$ & SPICE transmission line model (Bode plots) \\
\bottomrule
\end{tabular}
\end{table}

\section{Summary: Zero-Parameter Count}

The engine contains:
\begin{itemize}
    \item \textbf{0} trainable / fitted parameters
    \item \textbf{0} empirical structural data (PDB coordinates enter only as comparison targets)
    \item \textbf{12} derived physical constants (Table~\ref{tab:engine_constants_protein}), all traceable to Axioms~1--4
    \item \textbf{20} complex $Z_{topo}$ values (Table~\ref{tab:z_topo_complex}), derived from the SPICE backbone impedance model
    \item \textbf{1} objective function: $\overline{|S_{11}|^2}$
    \item \textbf{8} physics layers composing the loss, each with a clear electromagnetic interpretation
\end{itemize}

The force on each residue $\mathbf{F}_i = -\partial \mathcal{L} / \partial \mathbf{r}_i$ is computed exactly by reverse-mode automatic differentiation in a single forward--backward pass.  No finite-difference approximation, no force constants, no statistical potentials.
